% Figure 15.2 : ce qui se passe après kubectl create deployment (événements réels du chapitre 15).
\documentclass[tikz,border=4pt]{standalone}
\usepackage{quai}
\begin{document}
\begin{tikzpicture}[x=1cm,y=1cm,
    acteur/.style={qbox, font=\ttfamily\scriptsize, minimum width=2.4cm, minimum height=0.8cm, inner sep=2pt, align=center},
    m/.style={qarr},
    mt/.style={font=\scriptsize, fill=QPaper, inner sep=1.2pt}]

  \foreach \x/\n/\t/\c in {0/kc/kubectl/QGris, 3/api/API server\\+ etcd/QBleu, 6/dc/contrôleur\\Deployment/QSarcelle, 9/rc/contrôleur\\ReplicaSet/QSarcelle, 12/sc/scheduler/QSarcelle, 15/kl/kubelet\\+ containerd/QVert} {
    \node[acteur, draw=\c, fill=\c Bg] (\n) at (\x,0) {\t};
    \draw[QMuted, line width=0.5pt, dashed] (\x,-0.45) -- (\x,-8.6);
  }
  \newcommand{\msg}[5]{%
    \draw[m] (#1,#3) -- node[mt, above] {#4} (#2,#3);
    \node[qnum=QBleu] at ({#1+0.45*sign(#2-#1)},#3-0.32) {\scriptsize #5};}
  \msg{0}{3}{-1.0}{crée un Deployment}{1}
  \msg{3}{6}{-1.9}{nouveau Deployment}{2}
  \msg{6}{3}{-2.8}{crée un ReplicaSet}{3}
  \msg{3}{9}{-3.7}{nouveau ReplicaSet}{4}
  \msg{9}{3}{-4.6}{crée un Pod, sans nœud}{5}
  \msg{3}{12}{-5.5}{Pod à placer}{6}
  \msg{12}{3}{-6.4}{nodeName = minikube}{7}
  \msg{3}{15}{-7.3}{Pod pour ce nœud}{8}
  \msg{15}{3}{-8.2}{status : Running}{9}
  \node[qlbl, font=\scriptsize, anchor=west, align=left] at (15.25,-7.8) {Pulled, Created,\\Started};
\end{tikzpicture}
\end{document}
